\chapter{The SystemVerilog Verification Toolkit}
\label{ch:toolkit}

\chapabstract{The language features that exist only for verification: dynamic
data types, classes, dynamic threads, constrained randomisation, functional
coverage and assertions. None of it is synthesisable and almost none of it has
a \VHDL{} counterpart. This chapter is the vocabulary, taught one construct at
a time and with no testbench architecture in the way; \cref{ch:tb} is where it
gets assembled into a working testbench.}

Everything in this chapter is a tool, and a tool is easier to understand once
you know what it is for. So, briefly, here is the shape of the thing being
built in \cref{ch:tb}.

A modern testbench does not wiggle pins. It creates \emph{transactions},
objects that describe one complete operation on the interface. It
\emph{randomises} them inside constraints that describe what the design is
allowed to be given. It hands them to a \emph{driver} that knows the pin-level
protocol, watches the pins with an independent \emph{monitor}, checks the
results in a \emph{scoreboard} against a reference model, and measures what
the run actually exercised with \emph{functional coverage}. Around all of it,
\emph{assertions} state protocol rules that must hold on every cycle whether
anyone is looking or not.

The sections follow the path a transaction takes through that machine rather
than the order of the \LRM{}: storage first, because everything else is built
out of queues; then the class that a transaction is; then the threads and the
mailbox that move it between components, because a driver and a monitor are
useless until they can run at once and talk to each other; then the solver
that fills it in, the coverage that measures what the solver produced, and the
assertions that check the pins underneath all of it. \code{program} blocks
come last, as history.

Read it once for the shape, then use it as a reference. Nothing here is worth
memorising in isolation, and nothing here is worth writing until
\cref{ch:tb} gives it a job.

\begin{notebox}{What this chapter deliberately does not do}
It does not build a testbench. Every construct is shown on the smallest
example that demonstrates it, and questions of the form "where should this
component live" or "who should own this mailbox" are left to
\cref{ch:tb}, which answers them all at once with a worked example. If a
section feels like it stops one step short of usefulness, that is the step
\cref{ch:tb} takes.
\end{notebox}


% ===========================================================================
\section{Dynamic data types}
\label{sec:beyond:dyn}
\label{sec:toolkit:dyn}
\index{queue}\index{associative array}\index{dynamic array}

A testbench has to remember things: what it sent and has not yet checked,
which addresses it has touched, how many of each kind of transaction it has
produced. The number of things to remember is not known when the simulation
elaborates, because it depends on how the run goes.

That is the problem \VHDL{} solves badly. \VHDL{} gives you fixed arrays,
access types with manual \code{new} and \code{deallocate}, and that is all;
building a queue means writing a linked list, and building a sparse memory
model means writing a hash table. \SV{} has three dynamic containers and a
string type built into the language, and they are the foundation everything
later in this chapter stands on.

\textbf{Dynamic arrays} have a size fixed at run time:

\begin{svcode}
int    da [];
da = new[16];          // allocate
da = new[32](da);      // grow, preserving the first 16 elements
$display("%0d", da.size());
da.delete();           // free
\end{svcode}

\textbf{Queues} are the workhorse. They are ordered, indexable, and grow at
either end:

\begin{svcode}
int q [$];                    // unbounded queue
int b [$:15];                 // bounded to 16 entries

q.push_back(3); q.push_front(1);
int head = q.pop_front();
q.insert(1, 99);              // insert 99 at index 1
q.delete(0);                  // remove element 0
$display("%0d entries, last = %0d", q.size(), q[$]);
int slice [$] = q[1:3];       // sub-queue
q = {};                       // empty it
\end{svcode}

A queue is what a \VHDL{} designer would build from an access type and a lot
of pointer code. Push at one end, pop at the other, and you have the
expected-value list that every scoreboard is made of.

\textbf{Associative arrays} are sparse maps with an arbitrary index type:

\begin{svcode}
logic [31:0] mem [logic [31:0]];   // sparse memory model
int          hits [string];        // counters keyed by name
byte         anyidx [*];           // wildcard index type

mem[32'h8000_0000] = 32'hCAFE_BABE;
if (mem.exists(addr)) rdata = mem[addr];
hits["write"]++;

logic [31:0] k;
if (mem.first(k)) do $display("%h : %h", k, mem[k]); while (mem.next(k));
$display("%0d locations touched", mem.num());
\end{svcode}

An associative array indexed by a 32-bit address models a 4\,GB memory that
costs one entry per location actually written. Doing this in \VHDL{} means a
hash table you write yourself. Note \code{exists()}, and use it. Reading a key
that was never written is not an error: it returns the default value of the
element type, so a missing entry is indistinguishable from an entry holding
zero. Some tools also \emph{create} the entry on that read
(\tool{Verilator}~5.020 does; \code{m.num()} goes up and \code{m.exists()}
starts returning 1), so a memory model that polls addresses it has not written
can quietly grow. Test with \code{exists()} first, always.

\textbf{Strings} are a real type, not an array of characters:

\begin{svcode}
string s = "cordic";
$display("%0d %s %s", s.len(), s.toupper(), s.substr(0,2));
s = {s, "_rot"};                       // concatenation
string msg = $sformatf("x=%0d y=%0d", x, y);
int    n   = s.atoi();
\end{svcode}

\code{\$sformatf} is the one to remember. It formats into a \code{string}
instead of printing, which is how every object in this chapter builds the
description of itself that a failure message will quote.

\subsection{Array manipulation methods}
\index{array!methods}

Every array-like object, fixed or dynamic or queue or associative, supports a
family of methods that take an optional \code{with} clause. The \code{with}
clause is an expression evaluated for each element, where \code{item} names
the element. \refdesignbook lists the ones you will use.

\begin{table}[htbp]
  \centering
  \caption{Array manipulation methods. \code{item} is the implicit iterator
  name in the \code{with} clause.}
  \label{tab:beyond:arraymethods}
  \small
  \begin{tabular}{@{}lll@{}}
    \toprule
    \textbf{Method} & \textbf{Returns} & \textbf{Example} \\
    \midrule
    \code{sum} & scalar & \code{q.sum()} \\
    \code{product} & scalar & \code{q.product()} \\
    \code{and}, \code{or}, \code{xor} & scalar & \code{q.xor()} \\
    \code{min}, \code{max} & queue of 1 & \code{q.max()[0]} \\
    \code{find} & queue & \code{q.find(x) with (x > 10)} \\
    \code{find\_first}, \code{find\_last} & queue of 0 or 1 & \code{q.find\_first with (item == k)} \\
    \code{find\_index} & queue of indices & \code{q.find\_index with (item < 0)} \\
    \code{unique} & queue & \code{q.unique()} \\
    \code{unique\_index} & queue of indices & \code{q.unique\_index()} \\
    \code{sort}, \code{rsort} & in place & \code{t.sort with (item.addr)} \\
    \code{reverse}, \code{shuffle} & in place & \code{q.shuffle()} \\
    \code{size}, \code{delete} & nothing & \code{q.delete()} \\
    \bottomrule
  \end{tabular}
\end{table}

Two idioms are worth memorising:

\begin{svcode}
// how many transactions were writes?
int nwr = trs.sum() with (int'(item.is_write));

// the index of the first entry matching an address, or an empty queue
int idx [$] = trs.find_index with (item.addr == a);
if (idx.size() == 0) $error("no match for %h", a);
\end{svcode}

\code{sum() with} returning a count is the standard counting idiom: the
\code{with} expression is evaluated per element and the results are added.
Note the cast to \code{int}, which is not decoration.

\begin{gotcha}{The reduction is defined to happen in the element's width}
1800-2017 says the accumulation for \code{sum} is performed in the type of the
array element, so for \code{bit b [] = '\{1,1,1,1\};} the additions are
one-bit-wide and \code{b.sum()} is defined to be \code{0}, the parity, not
\code{4}. Nothing warns you.

Worse, tools disagree. \tool{Verilator}~5.020 returns \code{4} for that
expression, so code that is silently wrong on a commercial simulator is
silently right on the open-source one, and the bug appears only when you
change tool. \tool{Icarus Verilog}~12.0 does not implement \code{sum} on a
dynamic array at all and refuses to elaborate. The portable form is the only
form worth writing: \code{b.sum() with (int'(item))}, every time you are
counting rather than adding.
\end{gotcha}

% ===========================================================================
\section{Classes and object orientation}
\label{sec:beyond:class}
\label{sec:toolkit:class}
\index{class}

A queue holds things. The next question is what those things are, and for a
testbench the answer is almost always a \emph{transaction}: one complete
operation on the interface, with the stimulus that went in, the result that
came back, and enough identity to find it again in a log of ten thousand.

A \SV{} class is what a \VHDL{} protected type almost is, plus inheritance,
plus a garbage collector. If you have used a protected type to build a shared
counter or a mutex, you already have the idea: data plus methods, with the
data reachable only through the methods. What \VHDL{} does not give you is a
protected type that can be \emph{extended}, copied, stored in a queue, or
randomised.

\begin{svcode}[caption={A transaction class, the basic unit of a \SV{} testbench},label={lst:beyond:class}]
class bus_txn;
  rand  bit [31:0] addr;
  rand  bit [31:0] data;
  rand  bit        is_write;
  local int        id;                 // not visible outside
  static int       next_id = 0;        // one copy for the whole class

  function new(bit [31:0] a = 0);
    addr = a;
    id   = next_id++;
  endfunction

  virtual function string convert2string();
    return $sformatf("[%0d] %s @%h = %h", id,
                     is_write ? "WR" : "RD", addr, data);
  endfunction
endclass
\end{svcode}

The pieces, in \VHDL{} terms:

\begin{itemize}
  \item \code{new()} is the constructor. \code{bus\_txn t = new();} allocates
        an object and returns a \emph{handle} to it. There is no
        \code{deallocate}: objects are garbage collected when the last handle
        to them is dropped.
  \item \code{this} refers to the current object, exactly as in C++ or Java,
        and is needed only when a local name shadows a member.
  \item \kw{local} and \kw{protected} restrict visibility; the default is
        public, unlike a protected type where everything is private.
  \item \kw{static} members belong to the class rather than to any object, so
        \code{next\_id} counts every transaction ever created.
  \item \kw{virtual} on a method enables polymorphism.
\end{itemize}

Two properties of a class matter more than the syntax. The first is that the
locals of a class method have automatic lifetime, which is what makes it safe
to call the same method from several threads at once
(\cref{sec:toolkit:lifetime}). The second is that a class variable is a
reference, and that one catches everybody.

\subsection{Handles are references}

This is the single most common source of confusion for a \VHDL{} designer, so
it gets its own box.

\begin{gotcha}{A class variable is a handle, and assignment copies the handle}
\begin{svcode}
bus_txn a = new();
a.addr = 32'h1000;
bus_txn b = a;      // b and a name the SAME object
b.addr = 32'h2000;
$display("%h", a.addr);   // 2000, not 1000
\end{svcode}
Assigning one handle to another does not copy the object; it makes a second
name for it. A \VHDL{} record assignment copies the whole record, and that
instinct is wrong here. The symptom is always the same: a container that
appears to hold $N$ different objects and actually holds $N$ references to one
object, which carries whatever values were last written into it. Put a
\emph{copy} into the container, or construct a fresh object each time round
the loop, and never store the handle you are about to reuse.
\end{gotcha}

Copying is explicit, and comes in two depths:

\begin{svcode}
bus_txn c = new a;          // shallow copy: new object, members copied
                            // ... but nested handles still shared

function bus_txn clone();   // deep copy: written by hand
  clone = new();
  clone.addr = addr; clone.data = data; clone.is_write = is_write;
  clone.payload = payload.clone();   // recursive, for nested objects
endfunction
\end{svcode}

\code{new a} performs a shallow copy: scalar members are duplicated, but a
member that is itself a class handle is copied as a handle, so both objects
share the nested object. UVM automates this with field macros
(\cref{ch:uvm}); without UVM you write \code{clone()} yourself.

\subsection{Inheritance and polymorphism}
\index{class!inheritance}\index{polymorphism}

\begin{svcode}
virtual class txn_base;                 // abstract: cannot be instantiated
  pure virtual function string convert2string();
  virtual function void display();
    $display("%s", convert2string());   // works for every subclass
  endfunction
endclass

class mem_txn extends txn_base;
  rand bit [31:0] addr;
  function new(); super.new(); endfunction
  virtual function string convert2string();
    return $sformatf("addr=%h", addr);
  endfunction
endclass

txn_base t;                             // a base handle ...
mem_txn  d = new();
t = d;                                  // ... holding a derived object
t.display();                            // calls mem_txn's convert2string
\end{svcode}

A \kw{virtual} method is resolved by the type of the \emph{object}, not the
type of the handle. A \kw{pure virtual} method has no body and forces every
subclass to supply one; a class containing one must be declared
\kw{virtual class} and cannot be instantiated. This is the mechanism the whole
of UVM is built on: the framework holds base-class handles and calls virtual
methods, and your subclasses supply the behaviour.

Going the other way, from a base handle to a derived one, needs a checked
cast:

\begin{svcode}
mem_txn d2;
if (!$cast(d2, t)) $fatal(1, "handle is not a mem_txn");
\end{svcode}

\code{\$cast} returns 0 rather than failing when the object is not of the
target type, so its return value must be checked.

\subsection{Parameterised classes and the scope operator}

\begin{svcode}
class fifo_model #(type T = int, int DEPTH = 8);
  local T q [$];
  static int instances = 0;
  function void put(T x);
    if (q.size() >= DEPTH) $error("overflow");
    q.push_back(x);
  endfunction
endclass

fifo_model #(bus_txn, 16) f = new();
$display("%0d", fifo_model#(bus_txn,16)::instances);   // :: reaches statics
\end{svcode}

The \code{::} operator reaches into a class scope without an object: static
members, static methods, typedefs and parameters. Note that
\code{fifo\_model\#(int,8)} and \code{fifo\_model\#(bus\_txn,16)} are
different types with separate static variables, which is a surprise the first
time a counter you believed was global turns out to have been counting two
separate populations.

\begin{rosetta}
\begin{rleft}
type counter_t is protected
  procedure incr;
  impure function get return natural;
end protected;

type counter_t is protected body
  variable n : natural := 0;
  procedure incr is begin
    n := n + 1;
  end procedure;
  impure function get return natural
  is begin return n; end function;
end protected body;
\end{rleft}
\begin{rright}
class counter_t;
  local int n = 0;
  function void incr(); n++; endfunction
  function int get(); return n; endfunction
endclass

// ... plus inheritance, copying,
// randomisation, queues of them,
// and no explicit deallocation.
counter_t c = new();
\end{rright}
\end{rosetta}

Say it plainly: \textbf{classes are not synthesisable}. Not "poorly
supported"; not synthesisable at all, by any tool, by design. They exist for
testbenches, and every line of class code you write is simulation-only.

% ===========================================================================
\section{Concurrency control}
\label{sec:beyond:fork}
\label{sec:toolkit:fork}
\index{fork!join}\index{mailbox}\index{semaphore}

Objects on their own do nothing. Two of the components in \cref{ch:tb}, the
driver and the monitor, have to run at the same time, at different rates, for
the whole simulation, and they have to pass transaction objects between
themselves without either one waiting on the other's schedule. That needs
threads that can be created after elaboration, and a way for threads to hand
each other data safely.

\VHDL{} has neither. A \VHDL{} process is static: it exists from elaboration
to the end of simulation, and the number of processes is fixed when the design
is elaborated. If you want ten concurrent bus agents you write ten processes,
or one generate loop that writes them for you, and you decide how many before
the run starts.

\SV{} threads are dynamic. You can create one at run time, wait for some of
them, kill the rest, and create more.

\subsection{\texttt{fork} and its endings}

\begin{svcode}[caption={The three flavours of fork},label={lst:beyond:fork}]
fork                      // join: wait for ALL branches
  drive_a();
  drive_b();
join

fork                      // join_any: continue after the FIRST finishes
  begin wait_for_response(); resp_ok = 1; end
  begin #1000ns;           resp_ok = 0; end   // a timeout
join_any
disable fork;             // kill whatever is still running

fork                      // join_none: continue immediately
  monitor_forever();      // a background thread
join_none

// later:
wait fork;                // block until all my child threads are done
\end{svcode}

\code{join} is a barrier and the one ending with a \VHDL{}-like feel. The
\code{join\_any} plus \code{disable fork} pair in the middle is the canonical
timeout idiom, and there is no way to write it in \VHDL{} without a shared
variable and a lot of care. \code{join\_none} starts a background thread and
returns immediately, which is how a component that must run forever gets
started without blocking whoever started it.

\code{wait fork} is the fourth ending and the one people forget. It is not a
\code{fork} at all: it is a statement that blocks until every thread the
current thread has spawned, by any of the three forms, has finished.

\begin{gotcha}{\texttt{disable fork} is scoped to the thread, not to the \texttt{fork}}
This is a language trap rather than a testbench trap, and it is worth
understanding precisely because the damage is silent.

\code{disable fork} terminates \emph{all} active descendant threads of the
thread that executes it. Not the branches of the immediately preceding
\code{fork}: every descendant, however long ago it was started and by which
call. If a task you invoked twenty lines earlier left a \code{join\_none}
thread running in the background, that thread dies too, without a message.

The fix is to give the \code{disable} a bounded scope by nesting the timeout
inside its own thread, so that the only descendants in existence are the two
branches you care about:

\begin{svcode}
fork begin : timeout_scope
  fork
    wait_for_response();
    #1us;
  join_any
  disable fork;      // sees only timeout_scope's own children
end join
\end{svcode}

The alternative is \code{disable} with an explicit block label, which kills
exactly that named block and nothing else. Either is fine; a bare
\code{disable fork} in a task that anybody else may call is not.
\end{gotcha}

\subsection{Passing data between threads}

\SV{} provides two library classes and one language construct for
inter-thread communication. All three are built into the language; none of
them has to be imported.

A \textbf{mailbox} is a typed, blocking FIFO. It is what a \VHDL{} designer
would build from a shared variable and a protected type, except that the
blocking semantics are already correct and already race-free.

\begin{svcode}
mailbox #(bus_txn) mbx  = new();    // unbounded: put() never blocks
mailbox #(bus_txn) mbx4 = new(4);   // bounded to 4 entries

mbx.put(t);              // blocks while the mailbox is full
mbx.get(t);              // blocks while it is empty, then removes
mbx.peek(t);             // like get, but leaves the item in place
if (mbx.try_put(t)) ...  // non-blocking: returns 0 instead of waiting
if (mbx.try_get(t)) ...
if (mbx.try_peek(t)) ...
$display("%0d pending", mbx.num());
\end{svcode}

Bounded and unbounded are genuinely different objects, not a tuning knob. An
unbounded mailbox is a pure queue with a blocking read: the producer runs as
fast as it likes and memory absorbs the difference. A bounded one is a
rendezvous with slack $n$: the producer is throttled to the consumer's rate,
which is usually what you want and is always what makes a stalled consumer
visible.

Note that a mailbox is untyped unless you parameterise it. Plain
\code{mailbox mbx = new();} accepts anything, and a type error then becomes a
run-time failure inside \code{get} rather than a compile error. Write
\code{mailbox \#(T)}.

A \textbf{semaphore} is a counting mutex. It holds a number of keys; a thread
that wants the resource takes keys and blocks until enough are available, and
returns them when it is done.

\begin{svcode}
semaphore sem = new(1);    // one key: a plain mutex
sem.get(1);                // blocks until a key is free
// ... exclusive section ...
sem.put(1);                // return it

if (sem.try_get(1)) ...    // non-blocking
\end{svcode}

\code{new(0)} creates a semaphore with no keys, which is a barrier that some
other thread must open with \code{put}. Note that \code{put} does not check
that you ever held a key: putting more keys back than you took quietly
increases the count and destroys the mutual exclusion.

A \textbf{named event} is a synchronisation point with no payload. It is
declared with \kw{event}, triggered with \code{->}, and waited on in one of
two ways that behave differently.

\begin{svcode}
event e;
-> e;                              // trigger it
@(e);                              // wait for the NEXT trigger
wait (e.triggered);                // ... or for one in this time step

wait_order(a, b, c);               // events must occur in this order
\end{svcode}

\begin{gotcha}{\texttt{@(e)} and \texttt{wait (e.triggered)} are not the same wait}
\code{@(e)} is edge-sensitive: it blocks until the \emph{next} trigger of
\code{e}, and a trigger that already happened is gone. \code{e.triggered} is a
state that is true for the whole time step in which the event fired, so a
thread that reaches \code{wait (e.triggered)} later in the same time step
still sees it.

The failure this causes is total and silent: a thread waits forever on an
event that fired a few picoseconds of simulated time before it got there, and
the testbench simply stops making progress with no error. It is at its most
likely exactly where you least want it, between a thread that was just started
with \code{join\_none} and has not yet executed its first statement, and the
thread that started it. Prefer \code{wait (e.triggered)} everywhere except in
a deliberate "wait for the next one" loop.
\end{gotcha}

\subsection{Static and automatic lifetime}
\label{sec:toolkit:lifetime}
\index{automatic lifetime}\index{static lifetime}

Dynamic threads bring a trap that does not exist in a language where the
number of processes is fixed. It catches everybody exactly once.

In Verilog, and by inheritance in \SV{}, the arguments and locals of a task or
function declared in a module, interface or package have \textbf{static}
lifetime by default: one copy, created at elaboration, shared by every
invocation. A \VHDL{} process variable is static in the same sense, so the
concept is not new. What is new is that \SV{} lets you call the same task from
several threads at once, and they then share those variables.

\begin{svcode}[caption={The bug: four threads, one variable},label={lst:toolkit:staticbug}]
task print_after(int n);      // `int n' is STATIC here
  #(n * 10ns);
  $display("thread %0d at %0t", n, $time);
endtask

initial begin
  for (int i = 1; i <= 4; i++) begin
    fork
      print_after(i);         // four concurrent calls, ONE copy of n
    join_none
  end
end
\end{svcode}

All four calls share the single static \code{n}. The first sets it to 1 and
suspends at the delay; the second sets it to 2 and suspends; by the time any
of them wakes, \code{n} is 4. The output is four identical lines, all at the
same time, and the testbench appears to have started one thread instead of
four.

There are two fixes and you need both.

\begin{svcode}[caption={The fix: automatic task, automatic loop variable},label={lst:toolkit:staticfix}]
task automatic print_after(int n);   // each call gets its own n
  #(n * 10ns);
  $display("thread %0d at %0t", n, $time);
endtask

initial begin
  for (int i = 1; i <= 4; i++) begin
    automatic int k = i;             // a fresh k per iteration
    fork
      print_after(k);
    join_none
  end
end
\end{svcode}

\code{task automatic} gives every invocation its own copy of the arguments and
locals. \code{automatic int k = i;} inside the loop body gives every
\emph{iteration} its own variable to capture the loop counter into; without
it, the forked thread reads \code{i} after the loop has advanced and you get
the same symptom for a different reason.

\begin{gotcha}{Where the default flips, and why it bites on a refactor}
Everything declared inside a \kw{class} or a \kw{program} already has
automatic lifetime. Everything declared directly in a \kw{module},
\kw{interface} or \kw{package} does not. So the trap does not appear while you
are writing class-based code, and appears the moment you lift a working class
method out into a module-level task to share it, at which point the same
source text means something different. Write \kw{automatic} on every task and
function that is not deliberately keeping state between calls, or declare the
whole enclosing unit automatic (\code{module automatic tb;},
\code{program automatic test;}), which flips the default for everything inside
it. \VHDL{} procedures are always automatic, so nothing in your background
prepares you for this.
\end{gotcha}

% ===========================================================================
\section{Randomisation and constraints}
\label{sec:beyond:random}
\label{sec:toolkit:random}
\index{randomisation}\index{constraint}

The transaction class exists; something has to fill it in. A \VHDL{}
testbench describes stimulus \emph{procedurally}: a loop that walks addresses,
a nested loop that walks data patterns, a hand-written list of corner cases.
You write the sequence of inputs, so the test contains exactly as many cases
as you had patience for.

\SV{} lets you describe the \emph{legal input space} declaratively and asks a
solver for points inside it. This is a different way of thinking, and it is
the single largest productivity difference between the two languages in
verification: the test now contains one description and a number you can
raise from a hundred to a hundred thousand overnight.

\begin{svcode}[caption={A constrained transaction},label={lst:beyond:rand}]
class bus_txn;
  rand  bit [31:0] addr;
  rand  bit [7:0]  len;
  rand  bit [1:0]  size;        // weighted below; see the note on randc
  rand  bit        is_write;

  constraint c_align  { addr[1:0] == 2'b00; }
  constraint c_range  { addr inside {[32'h0000_0000 : 32'h0000_FFFF],
                                     [32'h8000_0000 : 32'h8000_0FFF]}; }
  constraint c_len    { len inside {[1:16]};
                        is_write -> len <= 8; }
  constraint c_dist   { size dist {0 := 1, 1 := 2, 2 := 5, 3 := 2}; }
  constraint c_noxing { (addr[11:0] + len*4) < 4096; }  // no 4k crossing
endclass

bus_txn t = new();
if (!t.randomize()) $fatal(1, "constraints unsatisfiable");
\end{svcode}

Every variable marked \kw{rand} is chosen by the solver so that all active
constraints hold simultaneously. The solver considers the constraints as a
system, not in order: \code{c\_noxing} can force \code{len} down even though
\code{len} appears first in the file. That is the property that makes the
declarative form worth learning, and it is also why the distribution of the
result is rarely the one you pictured.

\index{randc}\kw{randc} is the other declaration you will meet. It makes a
variable \emph{cyclic}: the solver visits every value of its range once before
any value repeats, which is how you sweep a small field without writing the
loop. It is tempting here, because \code{size} has only four values. It is not
allowed here: 1800-2017 forbids a \code{dist} constraint on a \kw{randc}
variable, and for a good reason, since a cyclic variable already has a fixed
distribution. Choose one, weighted or cyclic, per variable.

\begin{gotcha}{Ignoring the return value of \texttt{randomize()}}
\code{randomize()} returns \code{1} on success and \code{0} when the solver
could not satisfy the constraints. It does not throw, print or stop.

Written as a bare statement, \code{t.randomize();} compiles, discards the
result, and leaves the object holding whatever values it had before, which for
a freshly constructed object is all zeros. The test then runs happily, sending
the same transaction a thousand times, and reports a full pass having verified
one vector. \code{void'(t.randomize());} is worse than the bare call, not
better, because it documents in the source that you meant to ignore a
failure.

The only acceptable forms are
\code{if (!t.randomize()) \$fatal(1, "...");} or a macro that expands to it.
Every serious lint rule flags a discarded \code{randomize()} return value, and
every team turns that rule on after being bitten once. When it does fail, use
the tool's constraint debugger to find the contradiction rather than deleting
constraints until the solver stops complaining.
\end{gotcha}

\begin{table}[htbp]
  \centering
  \caption{Constraint constructs.}
  \label{tab:beyond:constraints}
  \small
  \begin{tabular}{@{}ll@{}}
    \toprule
    \textbf{Construct} & \textbf{Meaning} \\
    \midrule
    \code{a inside \{[1:9], 15\}} & set and range membership \\
    \code{a -> b} & implication: if \code{a} holds then \code{b} must \\
    \code{if (a) b; else c;} & the same thing, written imperatively \\
    \code{x dist \{1 := 4, 2 := 1\}} & weight \emph{per value}, ratio 4:1 \\
    \code{x dist \{[1:4] :/ 8\}} & weight 8 shared \emph{across} the range \\
    \code{solve a before b} & ordering hint; changes the distribution only \\
    \code{foreach (arr[i]) arr[i] < 10;} & one constraint per element \\
    \code{unique \{a, b, c\}} & all different \\
    \code{soft x == 0;} & default, overridable by a later constraint \\
    \code{x.size() inside \{[1:8]\}} & constrain a dynamic array's size \\
    \bottomrule
  \end{tabular}
\end{table}

\subsection{Controlling randomisation from the outside}

The class states what is \emph{legal}. A test states what it wants
\emph{this time}, and it does so without editing the class. That separation is
the reason the mechanism is worth its complexity.

Inline constraints add requirements for one call only:

\begin{svcode}
if (!t.randomize() with { addr inside {[32'h1000:32'h1FFF]};
                          is_write == 1'b1; })
  $fatal(1, "cannot build a write into the low window");
\end{svcode}

Individual variables and constraint blocks can be switched off:

\begin{svcode}
t.size.rand_mode(0);        // stop randomising size; keep its value
t.c_align.constraint_mode(0);  // allow unaligned addresses for this test
\end{svcode}

And two hooks run around every call:

\begin{svcode}
function void pre_randomize();  ... endfunction   // set up state
function void post_randomize();                   // fix up derived fields
  crc = compute_crc(addr, data);
endfunction
\end{svcode}

For variables that are not class members there is \code{std::randomize()}:

\begin{svcode}
int a, b;
if (!std::randomize(a, b) with { a < b; b < 100; }) $error("no solution");
\end{svcode}

And for cases where a distribution is all you need, without a solver:

\begin{svcode}
x = $urandom();                  // 32-bit unsigned
y = $urandom_range(15, 4);       // inclusive, either argument order
randcase                         // weighted branch selection
  3: send_read();
  1: send_write();
  1: send_idle();
endrandcase
\end{svcode}

\begin{gotcha}{The solver does not distribute the way you expect}
Given \code{rand bit [3:0] a; rand bit b; constraint c \{ b -> a == 0; \}}
you might expect \code{b} to be 1 half the time. It is not. The solver picks
uniformly over the \emph{solution space}: there are 16 solutions with
\code{b == 0} and exactly one with \code{b == 1}, so \code{b} is 1 about one
time in seventeen. \code{solve b before a;} restores the intuitive split by
telling the solver to choose \code{b} first. This is not a bug and it is not
tool-specific; it follows from the definition. Whenever a random variable
constrains the range of another, check the distribution with a covergroup
before trusting it.
\end{gotcha}

\begin{tipbox}{Reproducibility}
A random test that cannot be replayed is worthless. Fix the seed on the
command line (\code{vsim -sv\_seed 12345} or the equivalent), print it in the
banner of every run, and record it with the failure. Do not call
\code{\$urandom} from code whose execution order depends on the number of
threads, because the per-thread random stability guarantees of the \LRM{}
depend on objects being created in a deterministic order.
\end{tipbox}

\begin{notebox}{Constraints need a real solver, and not every tool has one}
Constrained randomisation is the feature that most sharply divides simulators.
\tool{QuestaSim} and the other commercial simulators ship a constraint solver
and implement the whole of clause 18. The open-source tools do not, and they
fail in two different and equally dangerous ways.
\tool{Icarus Verilog}~12.0 rejects a \code{constraint} block outright
(\code{sorry: Constraint declarations not supported}), which is at least
honest. \tool{Verilator}~5.020 compiles the class, accepts
\code{randomize()}, and emits \code{\%Warning-CONSTRAINTIGN: Constraint
ignored (unsupported)} once per constraint block, then randomises the
\kw{rand} variables with every constraint discarded. Compiled exactly as
written, \refdesignbook produces five such warnings and a first draw
of \code{len = 156} against \code{len inside \{[1:16]\}}. A warning you
filtered out of the log is the only thing standing between that and a test
which believes its stimulus is legal.
\refsimulationbook returns to what this means for an open-source flow.
\end{notebox}

% ===========================================================================
\section{Functional coverage}
\label{sec:beyond:cov}
\label{sec:toolkit:cov}
\index{coverage!functional}\index{covergroup}

Random stimulus raises an immediate question: what did it actually test? The
solver was told what is legal, not what is interesting, and it has no idea
which of the legal points matter. The answer is functional coverage, and the
distinction from code coverage is the first thing to get straight, because the
audience for this book has usually met only the second.

\textbf{Code coverage} is measured by the tool from the \RTL{} source with no
help from you: which lines executed, which branches were taken both ways,
which expressions took each value, which state machine states and transitions
were visited. It answers "did the stimulus exercise the code I wrote?" It
cannot tell you about a feature you forgot to implement, because there is no
code for it.

\textbf{Functional coverage} is written by you, from the specification. It
answers "did the stimulus exercise the situations the specification cares
about?" A design that ignores byte enables can have 100\% code coverage and
0\% coverage of the byte-enable cross, because you wrote a coverpoint for
byte enables and the tests never varied them.

You need both. Code coverage finds untested code; functional coverage finds
untested intent.

\begin{svcode}[caption={A covergroup for the bus transaction},label={lst:beyond:cg}]
covergroup bus_cg @(posedge clk);
  option.per_instance = 1;
  option.at_least     = 2;      // a bin counts as hit after 2 samples

  cp_kind : coverpoint kind {
    bins rd   = {READ};
    bins wr   = {WRITE};
    bins idle = {IDLE};
  }

  cp_len : coverpoint len {
    bins one       = {1};
    bins small     = {[2:7]};
    bins burst[4]  = {[8:63]};       // auto-split into 4 bins
    bins max       = {64};
    illegal_bins bad = {0};
    ignore_bins  rsv = {[65:255]};
  }

  cp_addr : coverpoint addr[15:12] {
    wildcard bins low  = {4'b00??};
    bins     other     = default;
  }

  cp_seq : coverpoint kind {
    bins rd_then_wr = (READ => WRITE);
    bins back2back  = (WRITE [*3]);
  }

  x_kind_len : cross cp_kind, cp_len {
    ignore_bins no_idle_len = binsof(cp_kind.idle);
    bins        long_writes = binsof(cp_kind.wr) &&
                              binsof(cp_len) intersect {[8:64]};
  }
endgroup

bus_cg cg = new();      // covergroups must be constructed
\end{svcode}

The vocabulary, briefly:

\begin{itemize}
  \item A \kw{coverpoint} is a variable or expression whose values you care
        about. Without explicit bins, the tool creates one automatically per
        value, up to \code{option.auto\_bin\_max} (64 by default).
  \item \kw{bins} name value sets. \code{bins b[4] = \{[8:63]\}} splits the
        range into four; \code{bins b[] = \{[8:63]\}} makes one bin per value.
  \item \kw{illegal\_bins} turn a hit into a runtime error; \kw{ignore\_bins}
        remove values from the denominator, which is how you exclude
        unreachable cases honestly.
  \item \code{default} catches everything not named, and does not count
        towards coverage. It is a debugging aid, not a way to reach 100\%.
  \item Transition bins use \code{=>}: \code{(READ => WRITE)} covers the
        value pair on consecutive samples, and \code{[*3]} repeats.
  \item \kw{cross} builds the product of two coverpoints;
        \code{binsof}/\code{intersect} filter the product down to the
        combinations that mean something.
  \item \code{option.per\_instance = 1} keeps a separate record per covergroup
        instance instead of merging them, which you almost always want when
        the same agent is instantiated several times.
\end{itemize}

Sampling is either automatic, through the event in the \code{covergroup}
header, or manual:

\begin{svcode}
covergroup txn_cg (ref bus_txn t);   // no event: sampled explicitly
  cp_addr : coverpoint t.addr { bins lo = {[0:'hFFFF]};
                                bins hi = {['h8000_0000:'h8000_0FFF]}; }
endgroup
...
tcg.sample();
$display("coverage so far: %0.2f%%", tcg.get_coverage());
\end{svcode}

Explicit \code{sample()} is the better habit, because it lets you sample a
complete, validated transaction rather than whatever the wires happened to
hold at an arbitrary clock edge. \cref{ch:tb} shows where in a testbench that
call belongs.

\begin{gotcha}{A coverpoint with no bins can report 100\% and mean nothing}
\code{coverpoint addr;} on a 32-bit address does not create four billion bins.
It creates \code{auto\_bin\_max} of them (64), each covering a huge range, and
hitting one address in each range reports full coverage. Worse, a coverpoint
on a single-bit signal that never changes will report 50\% and look like
progress. Always write explicit bins that correspond to something in the
specification. Auto-bins are for a first look, never for a coverage sign-off
number.
\end{gotcha}

\begin{gotcha}{A covergroup is a type; nothing is sampled until you construct it}
\code{covergroup bus\_cg; ... endgroup} declares a type, exactly as
\code{class} does. Until \code{cg = new();} has run there is no instance and
every \code{sample()} call does nothing useful. The coverage report then shows
0\% and looks like a stimulus problem rather than a missing line. A covergroup
declared inside a class may only be constructed in that class's constructor,
which at least makes the right place to put the line obvious.
\end{gotcha}

\begin{ruleofthumb}
Write the covergroup before the tests, from the specification. A coverage
model written after the tests measures the tests, not the specification.
\end{ruleofthumb}

Coverage numbers are produced in order to be closed, and closing them is a
loop rather than a measurement: run many seeds, merge the databases, look at
the holes, and decide for each one whether it needs a directed test, a
constraint change or an honest \code{ignore\_bins}. That loop is a methodology
question rather than a language question, and \cref{sec:tb:closure} walks
through it with real numbers.

\begin{notebox}{Covergroups are a commercial-simulator feature}
Unlike constraints, where \tool{Verilator} silently ignores what it cannot
do, covergroups are rejected outright. \tool{Verilator}~5.020 reports
\code{\%Error-UNSUPPORTED: covergroup}, along with a separate error for each
coverpoint and each bin specification, and stops. \tool{Icarus
Verilog}~12.0 does not parse the construct at all: a \code{covergroup} block
is a syntax error. Functional coverage in this form needs
\tool{QuestaSim} or another commercial simulator.
\refsimulationbook describes what \tool{Verilator} offers instead, which
is line and toggle coverage, a different and much weaker thing.
\end{notebox}

% ===========================================================================
\section{SystemVerilog Assertions}
\label{sec:beyond:sva}
\label{sec:toolkit:sva}
\index{assertion}\index{SVA}

Coverage says what happened. Something still has to say what was allowed to
happen, on every cycle, whether or not a scoreboard was watching.

A \VHDL{} \code{assert ... report ... severity} statement checks a condition
at one instant. Checking something that spans time, such as "grant follows
request within three cycles", means writing a process with a counter, and the
process is usually longer and buggier than the design it checks.

SVA is a temporal language embedded in \SV{} for exactly this. It is also the
one construct in this chapter that a design engineer writes as often as a
verification engineer, because an assertion inside a module is a statement of
the module's contract, and \code{bind} (\refdesignbook) lets you
attach one to a module you do not own.

\subsection{Immediate and concurrent}

An \emph{immediate} assertion is the \VHDL{} one, in procedural code:

\begin{svcode}
always_comb begin
  a_onehot: assert ($onehot(sel))
    else $error("sel = %b is not one-hot", sel);
end
\end{svcode}

A \emph{concurrent} assertion is a separate, clocked, always-running check:

\begin{svcode}
a_gnt: assert property (@(posedge clk) disable iff (!rst_n)
                        req |-> ##[1:3] gnt)
       else $error("no grant within 3 cycles of request");
\end{svcode}

\subsection{Sampling semantics, and why they matter}
\index{assertion!sampling}

A concurrent assertion does not read signals when it runs. It reads
\emph{sampled values}, taken in the Preponed region of the clock tick, before
any design process has executed in that time step. This is the same idea as
the clocking-block input skew of \refdesignbook, and it exists for
the same reason: an assertion must see the values the flip-flops saw, not the
values they are producing.

That gives you a set of sampled-value functions, all of which operate on the
preponed values:

\begin{svcode}
$past(sig)      // the sampled value one clock tick ago
$past(sig, 3)   // three ticks ago
$rose(sig)      // sampled value went 0 -> 1 at this tick
$fell(sig)      // 1 -> 0
$stable(sig)    // unchanged since the previous tick
$changed(sig)   // the negation of $stable
\end{svcode}

\begin{gotcha}{An assertion sees the \emph{old} value at the edge}
\begin{svcode}
always_ff @(posedge clk) count <= count + 1;
a_inc: assert property (@(posedge clk) count == $past(count) + 1);
\end{svcode}
This passes, and the reason surprises people. At the edge where \code{count}
becomes 5, the assertion samples \code{count} in the Preponed region, so it
reads 4, and \code{\$past(count)} reads 3. The assertion is comparing the two
\emph{previous} values, and it never sees the value being assigned at the
current edge. Every concurrent assertion is one delta behind the procedural
code around it. Write \code{\$past} expressions in terms of what the register
held \emph{before} this edge, and cross-check the first cycle after reset,
where \code{\$past} of an unknown is \code{X} and comparisons quietly fail.
\end{gotcha}

\subsection{Sequences and properties}

\begin{svcode}[caption={A valid/ready handshake checker},label={lst:beyond:sva}]
module handshake_chk #(parameter int W = 16) (
  input logic clk, rst_n,
  input logic valid, ready,
  input logic [W-1:0] data
);
  // 1. valid must not be withdrawn before the transfer completes
  property p_valid_stable;
    @(posedge clk) disable iff (!rst_n)
      valid && !ready |=> valid;
  endproperty

  // 2. the payload must not change while the transfer is pending
  property p_data_stable;
    @(posedge clk) disable iff (!rst_n)
      valid && !ready |=> $stable(data);
  endproperty

  // 3. every request is eventually granted, within 32 cycles
  property p_no_starve;
    @(posedge clk) disable iff (!rst_n)
      $rose(valid) |-> ##[1:32] (valid && ready);
  endproperty

  a_valid_stable: assert property (p_valid_stable)
                  else $error("valid dropped at %t", $time);
  a_data_stable : assert property (p_data_stable);
  a_no_starve   : assert property (p_no_starve);

  // measurement, not checking:
  c_transfer    : cover property (@(posedge clk) valid && ready);
  c_backpressure: cover property (@(posedge clk) valid && !ready [*4]);
endmodule
\end{svcode}

\refdesignbook is a complete, useful checker. It is a plain module, so
it can be instantiated normally, or attached to a design you do not own with
\code{bind} (\refdesignbook); \refexamplebook attaches this one to the
CORDIC rotator that way. Three assertions and two covers replace perhaps
eighty lines of \VHDL{} checking process.

The operators, in order of how often you will need them, appear in
\refdesignbook.

\begin{table}[htbp]
  \centering
  \caption{The assertion operators worth learning first.}
  \label{tab:beyond:sva}
  \small
  \begin{tabular}{@{}ll@{}}
    \toprule
    \textbf{Operator} & \textbf{Meaning} \\
    \midrule
    \code{a |-> b} & overlapping implication: \code{b} starts on the same tick \\
    \code{a |=> b} & non-overlapping: \code{b} starts on the next tick \\
    \code{\#\#n} & delay of exactly \code{n} clock ticks \\
    \code{\#\#[a:b]} & delay between \code{a} and \code{b} ticks \\
    \code{\#\#[1:\$]} & eventually, unbounded (needs a liveness-aware tool) \\
    \code{disable iff (c)} & abandon any attempt in progress while \code{c} \\
    \code{a [*n]} & \code{a} holds on \code{n} consecutive ticks \\
    \code{a [->n]} & \code{a} holds on the \code{n}-th occurrence, gaps allowed \\
    \code{a [=n]} & \code{n} occurrences, the last need not be at the end \\
    \code{a throughout b} & \code{a} holds for the whole duration of \code{b} \\
    \code{a within b} & sequence \code{a} occurs inside sequence \code{b} \\
    \code{a intersect b} & both hold and end on the same tick \\
    \code{first\_match(a)} & take only the earliest match of \code{a} \\
    \code{a and b}, \code{a or b} & both / either sequence matches \\
    \code{not p} & the property must never hold \\
    \bottomrule
  \end{tabular}
\end{table}

\begin{figure}[htbp]
  \centering
  \begin{tikzpicture}[
      >={Latex[length=1.6mm]},
      att/.style={draw=accent,line width=1pt,->},
      fail/.style={draw=warnframe,line width=1pt,->},
      tick/.style={draw=codeframe,line width=0.6pt},
      lbl/.style={font=\scriptsize\sffamily}]

    % time axis
    \draw[->,line width=0.8pt] (0,0) -- (11.6,0) node[right,lbl] {$t$};
    \foreach \i in {0,...,6} {
      \draw[tick] (\i*1.6,-0.15) -- (\i*1.6,3.0);
      \node[lbl,below] at (\i*1.6,-0.2) {$c_{\i}$};
    }

    % req waveform
    \node[lbl,anchor=east] at (-0.15,0.6) {\ttfamily req};
    \draw[line width=1pt,draw=codekey]
      (0,0.4) -- (1.6,0.4) -- (1.6,0.9) -- (4.8,0.9) -- (4.8,0.4) -- (9.6,0.4);

    % attempts
    \node[lbl,anchor=east] at (-0.15,1.6) {attempts};
    \draw[att] (0,1.4) -- (0.5,1.4) node[right,lbl] {vacuous};
    \draw[att] (1.6,1.9) -- (4.8,1.9)
        node[midway,above,lbl] {attempt at $c_1$: matches at $c_3$};
    \draw[fail] (3.2,2.5) -- (9.6,2.5)
        node[midway,above,lbl] {attempt at $c_2$: no grant, fails at $c_6$};
  \end{tikzpicture}
  \caption{A concurrent assertion starts a new attempt at every clock tick.
  Ticks where the antecedent is false pass vacuously; several attempts may be
  in flight at once, and each one succeeds or fails independently.}
  \label{fig:beyond:sva}
\end{figure}

\refdesignbook shows the mental model that most often goes missing. A
concurrent assertion is not a single check. At \emph{every} tick of its clock,
a new evaluation attempt begins. If the antecedent (the left side of
\code{|->}) is false, the attempt passes \emph{vacuously} and costs nothing.
If it is true, the attempt stays alive until the consequent succeeds or fails.
Several attempts overlap in time. This is why \code{first\_match} exists and
why a badly written property can report several failures for one bug.

\subsection{\texttt{assert}, \texttt{assume}, \texttt{cover}, \texttt{restrict}}

The same property can be used four ways, and the meaning depends on the tool.

\begin{itemize}
  \item \kw{assert}: in simulation, check and report on failure. In formal, a
        proof obligation.
  \item \kw{assume}: in simulation, usually checked like an assert but
        reported as an environment error. In formal, a \emph{constraint} on
        the inputs, so the prover only considers traces where it holds.
        Assuming too much is the classic way to prove a false theorem.
  \item \kw{cover}: never fails. It reports how many times the sequence
        occurred. In formal, a request to find a trace that reaches it, which
        is also the best way to sanity-check that your assumptions have not
        made the design unreachable.
  \item \kw{restrict}: formal only, ignored by simulators. A constraint used
        to limit the proof space deliberately, for example to prove a property
        only for aligned addresses.
\end{itemize}

\begin{rosetta}
\begin{rleft}
-- PSL, in a VHDL comment or,
-- since VHDL-2008, directly:
default clock is rising_edge(clk);

property p_gnt is
  always (req -> next_e[1 to 3] (gnt));

assert p_gnt
  report "no grant in 3 cycles";
\end{rleft}
\begin{rright}
property p_gnt;
  @(posedge clk)
    req |-> ##[1:3] gnt;
endproperty

a_gnt: assert property (p_gnt)
  else $error("no grant in 3");
\end{rright}
\end{rosetta}

If you have met PSL, SVA will feel familiar: both descend from the same
temporal-logic tradition, both have a boolean layer, a temporal layer and a
verification layer, and the operators map closely (\code{next\_e[1 to 3]}
against \code{\#\#[1:3]}, \code{until} against \code{throughout}). The
practical differences are that SVA is part of the language rather than a
separate standard, that its tool support is far wider, and that it can call
into \SV{} functions and covergroups directly. PSL remains a reasonable choice
for a \VHDL{} design, and \tool{QuestaSim} accepts it in both the VHDL and the
Verilog flavours.

\begin{notebox}{Where to put assertions}
Three places, in increasing order of preference for \RTL{}: inline in the
module (fine for small, obviously-local checks such as one-hot state);
in a separate checker module bound with \code{bind} (the default for anything
non-trivial); or in the interface itself, so that every instance of the
protocol is checked wherever it appears. The last is how professional
verification IP ships its protocol checks.
\end{notebox}

\begin{notebox}{Which half of SVA an open-source tool gives you}
SVA support is not all-or-nothing, and \tool{Verilator}~5.020 draws the line
in a place worth knowing. With \code{--assert} it accepts immediate
assertions, \code{assert property} and \code{cover property} with a clock and
\code{disable iff}, the implication operators \code{|->} and \code{|=>}, and
the sampled-value functions \code{\$past}, \code{\$rose}, \code{\$fell} and
\code{\$stable}. A property built from those runs, and a failure prints a
proper message naming the assertion label. What it rejects is the temporal
layer: any \code{\#\#} delay, including a plain \code{\#\#1}, gives
\code{\%Error-UNSUPPORTED: \#\# (in sequence expression)}, and any \code{[*n]}
repetition gives \code{\%Error-UNSUPPORTED: boolean abbrev}. So
\code{p\_valid\_stable} and \code{p\_data\_stable} from
\refdesignbook compile and run under \tool{Verilator}, and
\code{p\_no\_starve} and \code{c\_backpressure} do not. Write the
single-cycle properties first; they catch most protocol bugs and they are the
portable ones.
\end{notebox}

% ===========================================================================
\section{\texttt{program} blocks}
\label{sec:beyond:program}
\label{sec:toolkit:program}
\index{program block}

One construct remains, and it is here so that you recognise it rather than so
that you use it.

A \kw{program} block looks like a module but its code runs in the Reactive
region of the scheduler, after all design activity in that time step has
settled. It was introduced in \SV{}-3.1 to solve the testbench race problem:
a testbench written in an ordinary \code{initial} block competes with design
processes in the Active region, and whether it sees the old or the new value
of a signal depends on the simulator's scheduling.

\begin{svcode}
program automatic test (simple_bus.tb bus);
  initial begin
    bus.cb.req <= 1'b1;
    @(bus.cb);
    // ... implicit $finish when the program's initial blocks end
  end
endprogram
\end{svcode}

The honest modern view is that clocking blocks solved the same problem more
precisely, and better. Once every testbench signal goes through a clocking
block with defined input and output skews, the region a testbench statement
executes in stops mattering. UVM does not use program blocks; most current
methodology guides advise against them; \tool{Verilator} does not support
them. Recognise one when you meet it in older code, understand why it is
there, and do not write new ones. The one thing worth keeping from the design
is the automatic lifetime: \code{program automatic} was the original way to
flip the default of \cref{sec:toolkit:lifetime}, and \code{module automatic}
does the same job in code you would actually write today.

% ===========================================================================
\section{Where this goes next}
\label{sec:toolkit:next}
% ===========================================================================

You now have every construct a \SV{} testbench is made of, and no testbench.
That is deliberate, and it is the last thing this chapter has to say.

A queue, a class, a mailbox, a constraint block, a covergroup and a property
are individually unremarkable, and the reason to learn them is what they
become when they are put together: a component that reconstructs transactions
from pins without being told what was sent, a generator whose output you can
multiply by a thousand without writing a thousand cases, and a verdict that
can be trusted because it can be shown to fail. \cref{ch:tb} assembles exactly
these pieces, in this order, around the CORDIC rotator, and it also answers
the questions this chapter kept postponing: which component owns which
mailbox, where the covergroup is sampled, when the run is allowed to stop.
\cref{ch:uvm} then replaces the hand-written scaffolding with a standard
library, and changes nothing else.

Read \cref{ch:tb} next. If a construct here did not make sense in isolation,
it will make sense when you see the job it does.

% ===========================================================================
\section{Checklist}
\label{sec:toolkit:checklist}
% ===========================================================================

\begin{itemize}
  \item \textbf{A class handle is a reference.} \code{b = a} gives you two
        names for one object. Copy explicitly, or construct a fresh object
        each time round the loop.
  \item \textbf{Cast before you count.} \code{q.sum()} accumulates in the
        element's width, so write \code{q.sum() with (int'(item))} whenever
        the answer is a count. Tools disagree about this; the cast is right
        everywhere.
  \item \textbf{Declare every task and function} \kw{automatic}, or accept
        that two threads share one variable. Class and \kw{program} members
        already are; module, interface and package members are not.
  \item \textbf{Contain every \code{disable fork}} inside its own
        \code{fork ... join}, or use a labelled block. A bare one kills every
        descendant thread you have, including a monitor somebody else started.
  \item \textbf{Prefer \code{wait (e.triggered)} to \code{@(e)}.} The first is
        race-free within the time step; the second loses a trigger that has
        already happened and then waits forever.
  \item \textbf{Always check what \code{randomize()} returned.} A failed
        solve leaves the old values in place and says nothing.
        \code{void'(t.randomize())} is worse than the bare call.
  \item \textbf{Constrain the legal input space}, not the test you happen to
        want. The test belongs in an inline \code{with \{...\}} clause.
  \item \textbf{Seed deliberately and record the seed.} A random test you
        cannot reproduce is not a test.
  \item \textbf{Coverage answers a different question from checking.} Passing
        tells you nothing about what you exercised.
  \item \textbf{A covergroup with no \code{bins}} auto-generates them, and
        auto-generated bins on a 32-bit value are meaningless. Write them,
        and remember that the covergroup does nothing until \code{new()} runs.
  \item \textbf{Concurrent assertions sample in the Preponed region}, so they
        see the value that existed \emph{before} the edge. This is what you
        want, and it is not what you expect on the first day.
  \item \textbf{Check what your simulator implements.} Constraints,
        covergroups and the temporal layer of SVA are the three features an
        open-source tool is most likely to reject or, worse, to ignore.
\end{itemize}
